% Vietnamese translation for OI Public Library
% Source-File: IOI中国国家候选队论文/2003/王知昆/王知昆.doc
\documentclass[11pt]{article}
\usepackage{oipl-vn}

\title{Dùng tư tưởng cực đại hóa để giải bài toán hình chữ nhật con lớn nhất}
\author{Wang Zhikun}
\date{2003}

\begin{document}
\maketitle

\origtitle{Using extremal thinking for maximum subrectangle problems}
\sourcepath{IOI China candidate team papers / 2003 / Wang Zhikun / Wang Zhikun.doc}

\begin{abstract}
Bài viết giới thiệu cách dùng tư tưởng cực đại hóa cho một lớp bài toán gần
đây thường xuất hiện: tìm hình chữ nhật con lớn nhất, hình chữ nhật con tối
ưu, và các biến thể liên quan. Hai thuật toán có tính tổng quát được phân
tích, rồi thông qua một số ví dụ, bài viết bàn về kỹ thuật lựa chọn và sử
dụng chúng trong các mô hình khác nhau.
\end{abstract}

\noindent\textbf{Từ khóa:} hình chữ nhật, điểm chướng ngại, hình chữ nhật con
cực đại.

\section{Bài toán}

\term{Bài toán hình chữ nhật con lớn nhất} được mô tả như sau: trong một
lưới hình chữ nhật cho trước có một số điểm chướng ngại. Cần tìm hình chữ
nhật con lớn nhất nằm trong lưới, không chứa điểm chướng ngại nào ở bên trong,
và có các cạnh song song với trục tọa độ.

Một ví dụ tiêu biểu là bài \textit{Bãi tắm cho bò} trong Winter Camp 2002.
John muốn xây một bãi tắm hình chữ nhật trong nông trại hình chữ nhật; bãi
tắm không được chứa bất kỳ điểm cho sữa nào của bò trong phần trong, nhưng
điểm cho sữa có thể nằm trên biên. Với số điểm cho sữa \(S\le 5000\), kích
thước nông trại có thể tới \(30000\times 30000\). Phát biểu đầy đủ của bài
này nằm trong tệp phụ lục, đã được dịch riêng.

\section{Định nghĩa}

Gọi một hình chữ nhật con là \term{hợp lệ} nếu bên trong nó không chứa điểm
chướng ngại nào và các cạnh song song với trục tọa độ. Chướng ngại nằm trên
biên vẫn được chấp nhận.

Một hình chữ nhật con hợp lệ được gọi là \term{cực đại} nếu không tồn tại một
hình chữ nhật con hợp lệ lớn hơn chứa nó. Hình chữ nhật con hợp lệ có diện
tích lớn nhất trong tất cả các hình chữ nhật con hợp lệ được gọi là
\term{hình chữ nhật con lớn nhất}.

\paragraph{Định lý 1.}
Hình chữ nhật con lớn nhất trong một hình chữ nhật có điểm chướng ngại nhất
định là một hình chữ nhật con cực đại.

\paragraph{Chứng minh.}
Nếu hình chữ nhật lớn nhất \(A\) không cực đại, theo định nghĩa tồn tại một
hình chữ nhật hợp lệ lớn hơn chứa \(A\). Điều này mâu thuẫn với giả thiết
\(A\) là lớn nhất.

Định lý rất hiển nhiên, nhưng là điểm khởi đầu quan trọng: thay vì xét mọi
hình chữ nhật con, ta chỉ cần bảo đảm mọi hình chữ nhật con cực đại đều được
xét.

Trong phần dưới, giả sử toàn bộ hình chữ nhật có kích thước \(n\times m\), và
số điểm chướng ngại là \(S\).

\section{Thuật toán 1: liệt kê theo điểm chướng ngại}

\paragraph{Định lý 2.}
Một hình chữ nhật con hợp lệ là cực đại khi và chỉ khi mỗi cạnh của nó hoặc
chứa một điểm chướng ngại, hoặc trùng với biên của toàn bộ hình chữ nhật.

Nếu một cạnh không chứa chướng ngại và cũng không trùng với biên ngoài, ta có
thể đẩy cạnh đó ra ngoài một đoạn nhỏ mà vẫn không chứa chướng ngại trong
phần trong, nên hình chữ nhật chưa cực đại. Chiều ngược lại cũng trực tiếp:
khi bốn cạnh đều bị chặn bởi chướng ngại hoặc biên ngoài, hình chữ nhật không
thể mở rộng theo bất kỳ hướng nào.

Một cách làm thô là thêm bốn góc của hình chữ nhật lớn vào tập chướng ngại,
rồi liệt kê bốn biên của hình chữ nhật con theo các điểm mà chúng chứa. Sau
đó kiểm tra hình chữ nhật có hợp lệ hay không. Cách này có độ phức tạp
\(O(S^5)\), quá lớn.

Cải tiến đầu tiên là chỉ liệt kê hai biên trái và phải. Với hai biên đã cố
định, sắp xếp các điểm nằm giữa chúng theo tung độ; hai điểm liên tiếp theo
thứ tự này xác định một khoảng trống theo chiều dọc, từ đó sinh ra một hình
chữ nhật hợp lệ. Độ phức tạp giảm xuống \(O(S^3)\). Tuy nhiên nhiều hình chữ
nhật được sinh ra vẫn không cực đại, vì biên trái hoặc biên phải có thể không
bị chặn đúng cách.

Thuật toán 1 giảm tiếp lượng việc thừa bằng cách cố định biên trái rồi quét
các điểm phía bên phải theo thứ tự hoành độ tăng dần.

Với một điểm chướng ngại \(P\) được chọn làm điểm chặn biên trái, ban đầu đặt
biên trên và biên dưới hiện tại là hai biên ngoài của hình chữ nhật lớn. Sắp
xếp các điểm nằm bên phải \(P\) theo hoành độ rồi lần lượt quét:

\begin{itemize}
  \item Khi gặp một điểm \(Q\) nằm trong dải hiện tại, dùng hoành độ của \(Q\)
  làm biên phải. Cùng với biên trái và hai biên trên/dưới hiện tại, ta thu
  được một hình chữ nhật con cần xét.
  \item Nếu \(Q\) nằm phía dưới \(P\), biên dưới hiện tại phải được nâng lên
  tới tung độ của \(Q\), vì hình chữ nhật sau này không được chứa \(Q\).
  \item Nếu \(Q\) nằm phía trên \(P\), biên trên hiện tại được hạ xuống tung
  độ của \(Q\).
  \item Nếu \(Q\) cùng hàng với \(P\), quá trình quét có thể dừng vì các hình
  chữ nhật tiếp theo có chiều cao bằng \(0\).
\end{itemize}

Các điểm không nằm trong dải trên/dưới hiện tại không ảnh hưởng tới hình chữ
nhật đang xét và có thể bỏ qua. Vì ta đã thêm hai góc phải của hình chữ nhật
lớn vào tập điểm, trường hợp biên phải trùng với biên ngoài cũng không bị bỏ
sót.

Phần trên chỉ xét các hình chữ nhật có biên trái chứa một điểm chướng ngại.
Ta còn phải xử lý biên trái trùng với biên ngoài. Trường hợp biên phải chứa
một điểm chướng ngại được xử lý đối xứng bằng cách quét từ phải sang trái.
Trường hợp cả hai biên trái và phải đều trùng biên ngoài được xử lý trước:
sắp xếp mọi điểm theo tung độ, rồi xét các dải giữa hai tung độ liên tiếp.

Như vậy thuật toán 1 liệt kê được tất cả hình chữ nhật con cực đại. Độ phức
tạp thời gian là \(O(S^2)\). Thuật toán này rất hiệu quả khi số chướng ngại
thưa.

\section{Thuật toán 2: liệt kê dây treo}

Thuật toán 1 phụ thuộc vào số điểm chướng ngại. Nếu \(S\) có thể lớn tới
\(n m\), nó không còn đủ tốt. Ta cần một thuật toán phụ thuộc trực tiếp vào
diện tích lưới.

\paragraph{Định nghĩa.}
Một đoạn thẳng dọc được gọi là \term{đoạn dọc hợp lệ} nếu ngoài hai đầu mút,
nó không phủ lên điểm chướng ngại nào. Một \term{dây treo} là đoạn dọc hợp lệ
có đầu trên chạm một điểm chướng ngại hoặc chạm biên trên của hình chữ nhật
lớn.

Với mọi hình chữ nhật con cực đại, biên trên của nó hoặc chứa chướng ngại,
hoặc trùng biên trên của hình chữ nhật lớn. Nếu cắt hình chữ nhật con đó bằng
các đoạn dọc, bên trong nó luôn tồn tại một dây treo. Ngược lại, nếu lấy một
dây treo và trượt nó sang trái, sang phải xa nhất có thể mà vẫn hợp lệ, ta
thu được một hình chữ nhật con hợp lệ.

\paragraph{Định lý 3.}
Tập các hình chữ nhật con nhận được bằng cách trượt mọi dây treo sang trái và
sang phải xa nhất có thể chứa toàn bộ tập hình chữ nhật con cực đại.

Ở đây ``xa nhất'' nghĩa là trượt đến khi gặp chướng ngại hoặc biên ngoài.
Vì mỗi dây treo tương ứng với điểm đáy của nó, số dây treo là
\((n-1)m\). Nếu xử lý mỗi dây treo trong \(O(1)\), toàn bộ thuật toán có độ
phức tạp \(O(nm)\).

\subsection{Các công thức quy hoạch động}

Với dây treo có đáy tại ô hoặc điểm \((i,j)\), đặt:

\begin{itemize}
  \item \(\mathit{height}[i,j]\): chiều cao của dây treo;
  \item \(\mathit{left}[i,j]\): vị trí xa nhất có thể trượt sang trái;
  \item \(\mathit{right}[i,j]\): vị trí xa nhất có thể trượt sang phải.
\end{itemize}

Đồng thời, trên hàng \(i\), đặt \(L[i,j]\) là chướng ngại gần nhất bên trái
\((i,j)\), hoặc biên trái nếu không có; đặt \(R[i,j]\) là chướng ngại gần
nhất bên phải, hoặc biên phải nếu không có.

Nếu vị trí ngay phía trên \((i,j)\), tức \((i-1,j)\), là chướng ngại, dây
treo mới bắt đầu tại đây:
\[
  \mathit{height}[i,j]=1,\qquad
  \mathit{left}[i,j]=L[i,j],\qquad
  \mathit{right}[i,j]=R[i,j].
\]

Nếu \((i-1,j)\) không phải chướng ngại, dây treo tại \((i,j)\) là dây treo
tại \((i-1,j)\) kéo dài thêm một đơn vị:
\[
  \mathit{height}[i,j]=\mathit{height}[i-1,j]+1.
\]
Biên trái và phải phải đồng thời thỏa ràng buộc của dây cũ và của hàng mới:
\[
  \mathit{left}[i,j]=\max\{\mathit{left}[i-1,j],\,L[i,j]\},
\]
\[
  \mathit{right}[i,j]=\min\{\mathit{right}[i-1,j],\,R[i,j]\}.
\]

Diện tích hình chữ nhật do dây treo đáy \((i,j)\) sinh ra là
\[
  \bigl(\mathit{right}[i,j]-\mathit{left}[i,j]\bigr)
  \mathit{height}[i,j].
\]
Vì vậy đáp án là giá trị lớn nhất của biểu thức trên với mọi \((i,j)\). Mỗi
trạng thái được tính trong \(O(1)\), nên thuật toán 2 có thời gian và bộ nhớ
\(O(nm)\).

\section{So sánh hai thuật toán}

Hai thuật toán trên có phạm vi sử dụng khác nhau. Thuật toán 1 có độ phức tạp
\(O(S^2)\), phù hợp khi chướng ngại thưa. Thuật toán 2 có độ phức tạp
\(O(nm)\), không phụ thuộc trực tiếp vào số chướng ngại, nên phù hợp khi lưới
không quá lớn hoặc chướng ngại dày.

Trong một số bài, nếu kích thước tọa độ rất lớn nhưng số chướng ngại ít, có
thể rời rạc hóa tọa độ để dùng thuật toán 2. Tuy nhiên rời rạc hóa thường
làm cài đặt phức tạp hơn; khi \(S\) nhỏ, thuật toán 1 thường tự nhiên hơn.

\section{Ví dụ}

\subsection{\textit{Bãi tắm cho bò}}

Bài toán cho một hình chữ nhật và các điểm chướng ngại, yêu cầu tìm hình chữ
nhật hợp lệ có diện tích lớn nhất, đúng với mô hình đã phân tích.

Thuật toán 1 dùng trực tiếp, thời gian \(O(S^2)\), bộ nhớ \(O(S)\). Thuật toán
2 cần rời rạc hóa vì kích thước nông trại có thể tới \(30000\times 30000\);
sau rời rạc hóa, kích thước hiệu dụng khoảng \(S\times S\), thời gian và bộ
nhớ cũng là \(O(S^2)\) và \(O(S)\), nhưng cài đặt phức tạp hơn. Vì vậy với
bài này, thuật toán 1 tốt hơn cả về hiệu quả lẫn độ đơn giản.

\subsection{\textit{Candy}}

Bài này cho một hộp kẹo \(n\times m\). Ô \((i,j)\) có \(a[i,j]\) viên kẹo;
ô có \(a[i,j]=0\) bị xem là thủng. Cần cắt ra một hình chữ nhật không có lỗ
và có tổng số kẹo lớn nhất, với \(1\le n,m\le 1000\).

Đây là biến thể \term{hình chữ nhật con hợp lệ có trọng số lớn nhất}. Trong
ma trận, biên của hình chữ nhật cũng không được chứa ô thủng. Vì mọi trọng số
đều không âm, hình chữ nhật tối ưu vẫn phải là một hình chữ nhật cực đại:
nếu còn mở rộng được qua các ô hợp lệ, tổng kẹo không giảm.

Thuật toán 1 có thể tệ tới \(O(n^2m^2)\) vì số ô thủng có thể đạt \(nm\).
Thuật toán 2 chỉ cần sửa phần tính diện tích thành tính tổng trọng số bằng
tổng tiền tố hai chiều; độ phức tạp vẫn là \(O(nm)\). Vì vậy thuật toán 2 phù
hợp hơn.

\subsection{\textit{Big Barn}}

Bài \textit{Big Barn} yêu cầu đặt một hình vuông lớn nhất trong nông trại
\(N\times N\), tránh các điểm có cây, với \(N\le 1000\) và \(T\le 10000\).

Gọi hình vuông con hợp lệ là hình vuông không chứa chướng ngại trong phần
trong. Hình vuông con cực đại là hình vuông hợp lệ không thể mở rộng thêm.
Tương tự hình chữ nhật:

\paragraph{Định lý 4.}
Hình vuông con lớn nhất trong một hình chữ nhật có chướng ngại nhất định là
một hình vuông con cực đại.

Đối với hình vuông, điều kiện cực đại có dạng khác: một hình vuông con hợp lệ
là cực đại khi mọi cặp cạnh kề đều bị ít nhất một chướng ngại hoặc biên ngoài
chặn. Từ đó suy ra:

\paragraph{Định lý 5.}
Mỗi hình vuông con cực đại đều nằm trong ít nhất một hình chữ nhật con cực
đại, và có hai cạnh đối diện trùng với hai cạnh của hình chữ nhật con cực đại
đó.

Vì vậy chỉ cần liệt kê mọi hình chữ nhật con cực đại bằng một trong hai thuật
toán trên. Với mỗi hình chữ nhật kích thước \(a\times b\), hình vuông lớn
nhất mà nó chứa có cạnh
\[
  \min(a,b).
\]

Với \textit{Big Barn}, thuật toán 1 có độ phức tạp \(O(T^2)\), còn thuật toán
2 có độ phức tạp \(O(N^2)\). Do \(N\le 1000\) và \(T\) có thể tới \(10000\),
thuật toán 2 thường tốt hơn. Kết quả chạy thử trên USACO Training Gateway
trong bài gốc cũng cho thấy thuật toán 2 nhanh hơn rõ rệt trên dữ liệu lớn.

\section{Kết luận}

Thiết kế thuật toán nên bắt đầu từ đặc trưng cơ bản của bài toán. Với lớp bài
toán hình chữ nhật con lớn nhất, điểm đột phá là tư tưởng cực đại hóa: nghiệm
tối ưu nằm trong tập các nghiệm cực đại. Từ đó có thể xây dựng hai hướng
liệt kê: một hướng dựa vào điểm chướng ngại, phù hợp khi chướng ngại thưa; và
một hướng dựa vào dây treo trên toàn bộ lưới, phù hợp khi kích thước lưới
không quá lớn hoặc chướng ngại dày.

Hai thuật toán có vẻ rất khác nhau, nhưng bản chất giống nhau: đều dùng cực
đại hóa để giảm không gian tìm kiếm từ mọi hình chữ nhật con xuống một tập
ứng viên có cấu trúc. Khi giải bài thực tế, không nên chỉ máy móc áp dụng một
thuật toán có sẵn; cần phân tích mô hình, mật độ dữ liệu, giới hạn bộ nhớ và
độ phức tạp cài đặt để chọn biến thể phù hợp.

Phụ lục của bài gốc gồm phát biểu chi tiết và chương trình cho các ví dụ; nội
dung đó được dịch riêng trong bản phụ lục tương ứng, không lặp lại ở đây.

\begin{thebibliography}{9}
\bibitem{guide}
Wu Wenhu, Wang Jiande.
\textit{Informatics Olympiad Contest Guide: Analysis of 1997--1998 Problems}.

\bibitem{ioi99}
\textit{Selected Papers of the IOI 1999 China Training Team}.

\bibitem{quarterly}
\textit{Informatics Olympiad Quarterly}.

\bibitem{gold}
Jiang Wenzai, editor.
\textit{Road to Gold Medal: Contest Tutoring}.
Shaanxi Normal University Press.
\end{thebibliography}

\end{document}
